\subsection{ログ付き追加比較（小規模診断）}\label{sec5.additional}

保存結果の再集計とは別に，2026年9月6--7日に，状態変化の\texttt{test\_task1, test\_task6}と配置の\texttt{test\_task1, test\_task65}を用いた追加比較を行った．4タスク，2モデル，10条件，各3反復の計240エピソードである．モデルは\texttt{gpt-4o-mini-2024-07-18}と\texttt{gpt-4o-2024-08-06}，temperatureは0，各要求の出力上限は256トークンとした．初期位置を固定し，各試行のAPI呼出し前に参照グラフとの記号状態一致，全ノードの位置差0.001以下，四元数距離0.0001以下を確認した．全240件がこの照合を通過した．これは過去の人物位置を復元した実験ではない．

C0は前提条件検証なし，C1は既存のLLM検証，C2は同じ抽出知識と条件辞書による記号検証である．C3は検証せず実行し，VH失敗後の最新知識とエラーに基づき1回修復する独自対照で，失敗実行も参照列長の2倍の実行予算に含める．CAPE等の忠実な再現とは位置付けない．C5は検証器だけに完全グラフを与え，計画器への知識は変えない．同じ8操作の条件を評価するが，既知の空状態集合と情報不足の区別もC2と異なるため，差を抽出漏れだけに帰属させない．

C1-budgetとC4-budgetは，毎回の行動決定で生成・レビュー・説明・再生成の4要求を行う専用対照である．前者は条件辞書，後者は条件辞書を提示しない一般的再考を用いる．両者の出力上限合計は1024トークンだが，入力長，実際の出力長，計画全体の要求数は一致しない．元の必要時のみ再生成するC1とは別仕様であり，厳密な総計算量一致は主張しない．R-random，R-fixed，R-noneはC1の事例だけを，保存済みの反復別ランダム1例，固定1例，空の事例メッセージに変更する．C1が意味類似検索の基準となる．反復ごとにタスク・モデルのブロックと条件順をseed 20260907で並べ替え，事前に保存した．

\begin{table}[tb]
\caption{追加比較の成功数（各条件4タスク$\times$3反復，分母12）}\label{tab:step4_pilot}
\centering\small
\begin{tabular}{lrr}
\hline
条件 & GPT-4o mini & GPT-4o \\
\hline
C0 & 8 & 9 \\
C1 & 9 & 9 \\
C2 & 4 & 10 \\
C3 & 5 & 9 \\
C1-budget & 4 & 11 \\
C4-budget & 6 & 10 \\
C5 & 5 & 9 \\
R-random & 8 & 9 \\
R-fixed & 7 & 9 \\
R-none & 3 & 3 \\
\hline
\end{tabular}
\end{table}

\ref{tab:step4_pilot}には参照再生が失敗した配置1件も含めた．別途，参照成功3件のみの集計も保存した．初期充足の状態変化1件は全条件・全反復で行動せず成功し，表の各分子に3件ずつ含まれる．4タスクは3種類の指示名に対応し，12個の独立タスクを評価したわけではない．C1とC0，C1-budgetとC4-budgetなどの差の方向はモデルによって異なり，一律な優位性や全415タスクへの一般化を示すものではない．

判定過程について，C1の候補に対するLLM判定を，抽出辞書上のJSON条件のtrue/falseと照合した正答数は，miniで34/43（誤受理4，誤棄却5），4oで45/54（誤受理9，誤棄却0）だった．これらは異なる軌跡で観測された判定であり，共通の固定問題集合に対するモデル比較ではない．情報不足のラベルは分母から除く規則としたが，このC1集計には該当例がなかった．

全10条件を合わせた終了判定はminiで496回，4oで600回であり，未達なのに厳密な文字列Endを返した回数はそれぞれ26回，0回だった．4oにはEnd/Continue以外の応答が2回あり，書式違反として別記した．一方，完全グラフ上のJSON条件がtrueだった実行でも，miniは376回中28回，4oは447回中37回でVHが失敗を返した．条件辞書の充足はVHの全実行制約の充足を保証しない．棄却して実行していない候補にVH成否を付与したものではない．

完全グラフ，抽出知識，全プロンプト・応答，生成候補，検証，再生成，実行前後の状態を保存した．新ドライバはEnd判定の入力となった直前グラフで成否を計算し，終了後にも独立に最終グラフを取得する．元NotebookのEnd後の取得時点と厳密に同一ではないが，今回の240件では成功ラベルと最終ゴール充足の不一致はなかった．成功API要求は3047件，入力1052576トークン，出力30594トークンであり，通常単価・キャッシュ割引なしの推定費用は1.681766米ドルだった．通信失敗と出力上限による打切りは0件である．設定・ソースのハッシュ，初期状態，ゴール，入力・応答，費用を照合した．個別値，反復別・対応比較，費用と再実行手順は補足資料\texttt{Paper/vh\_step4/}に示す．
